% ---------------------------------------------------------------------------
% E4: held-out test of the decoded-word head sets (Sec. 15).
% Written to be pasted after Sec. 15 of proposal-simple.tex.
% Symbols: H^J_{t,s} as in the main text; A, B are the two halves of a task's items.
% ---------------------------------------------------------------------------

\section{E4: Held-out test of head sets from decoded words}
\label{sec:e4}

\subsection{Design}

\paragraph{Question.} In Sec.~15 the head sets chosen from the words heads decode
(CC-pooled) carry every task better than the E1 J-head sets, and carry tasks the E1 sets do
not (antonym, singular--plural, next-item, the three translations). That result is exploratory:
the words, the list sizes and the designs were chosen after seeing the first check on the same
items, each effect is the best of four values of $\alpha$, each set was run once, the second check
was not run, and the sets (95 to 137 heads) are larger than the E1 sets and than the function
vector of Davidson et al.\ (2025). E4 asks three questions with the design of Sec.~15 frozen:
(i) does CC-pooled still beat the E1 set on items not used to build it; (ii) is its effect larger
than that of random heads of the same number and layers; (iii) at the same number of heads, do
heads chosen by what they decode carry the task as well as heads chosen by their causal effect
(Davidson et al., 2025)? A second part, E4b (Sec.~\ref{sec:e4b}), adds new wordings of
each task's instruction and asks two more: (iv) do words and heads chosen on more wordings
carry the task better; (v) do heads chosen on some wordings carry the task when it is asked in
wordings never used to choose them? The hypotheses are stated before the run (below).

\paragraph{Frozen choices.} Everything that Sec.~15 chose after seeing results is fixed here
and not changed after E4 is run: J-head readout; rule style; word filter (alphabetic, at least
three letters, not a function word, Zipf frequency at least 3); task words by share on the task
minus the mean share on the other 13 tasks; common words by the largest minimum share over
the 14 tasks; 5 task words and 5 common words; heads by Definition~1 with $k = 10$,
$\rho = 0.10$, on the five templates' prompts pooled; for CC-pooled-clean the 50 words with the
largest share on the 1{,}000 WikiText-103 passages removed first; for attend and decode the
score of Sec.~15. No other list, list size, threshold or contrast is run in E4.

\paragraph{Item split.} The items of each task are shuffled with a fixed seed (0) and cut into
two halves, $A$ and $B$, of equal size (one more in $A$ when odd). Each half is used once to
choose and once to test (two folds): in fold~1 everything is chosen on $A$ and tested on $B$;
in fold~2 the reverse. ``Chosen'' covers every quantity that depends on the items: the
identification prompts $P^+$ of each template (the first 100 correct prompts of that half),
the word shares, the word lists, the head sets, the rankings of the top-$N$ sets, the
function-vector scores, the task means and the control means (the control task's own half
with the same name). ``Tested'' covers only the patched items. The WikiText readings do not
depend on the items and are shared by both folds. On tasks with fewer than 200 items a half
gives fewer than 100 identification prompts or patched items; the number used is reported for
every task, fold and check.

\paragraph{Head sets.} For each task and fold, all J-head:
\begin{enumerate}
  \item \textbf{E1}: Definition~1 with Yang et al.'s term list on the most accurate template
    of the choosing half (the E1 rule, recomputed on that half).
  \item \textbf{CC}: CC-pooled at its own size $N_t$ (the heads passing Definition~1).
  \item \textbf{CC-20, CC-40}: the first 20 and 40 heads of the model ranked by their pooled
    match rate under the CC list (ties as in Sec.~14); a head need not pass the threshold.
  \item \textbf{FV-$N$}, $N \in \{20, 40, N_t\}$: the first $N$ heads ranked by the
    function-vector score of Davidson et al.\ (2025) on task $t$ (the gain in the answer's
    probability when the head is set to its mean on prompts whose instruction is replaced, as
    run for the summary), computed on the choosing half only.
  \item \textbf{Random-$N$}, $N \in \{20, 40, N_t\}$: five draws per size, each with the same
    number of heads in every layer as the CC set of that size, drawn outside it (seeds 0--4).
  \item Secondary, not in the hypotheses: \textbf{CC-clean} (5\,+\,5 words, pooled), and
    \textbf{A+D-20, A+D-40} (attend and decode); \textbf{FV-univ-20, FV-univ-40}, the heads
    ranked by the mean of the function-vector score over the 14 tasks (Davidson et al.'s own
    choice of heads).
\end{enumerate}
The FV ranking uses a causal criterion on prompts without the instruction, close to the first
check itself; the comparison is therefore tilted towards the FV heads, not against them.

\paragraph{Causal checks.} Both checks of E1 (Sec.~5.2.1), on the testing half, with the
rule-style means of the five templates of the choosing half and the control mean of E1's
control task, at ten values of $\alpha$: the nine of Sec.~14 ($-10, -4, 2, 3, 4, 5, 6, 8, 10$),
which give within 0.05 of the fifteen on the E1 and E1b sets, and $\alpha = 1$, added because
the CC sets of Sec.~15 had their best effect at $\alpha \le 2$ on 11 of the 14 tasks. Check~1: the first 100 items of
the testing half answered correctly with one shot, patched in their zero-shot prompts. Check~2:
the first 100 answered wrongly with one shot, patched in their 1-shot prompts; unlike E1, the
control mean is run in check~2 as well. Random sets are patched with the task means only.
Scoring: the flexible score; the first-token score and the reply kinds of Sec.~10 are kept beside
it. As a reference, not a matched comparison, Davidson et al.'s own intervention (the 30-head
function vector of the summary, the sum of its heads' means added after layer 11 at the last
token, $\alpha = 1$) is run on the
same items.

\paragraph{Choice of $\alpha$.} Two effects are reported for every set, check and task. The
\emph{best-$\alpha$ effect}, as in the earlier sections: the highest accuracy over the ten
values, averaged over the five means, minus the unpatched accuracy. The \emph{cross-fitted
effect}, the primary measure: for fold~1, the accuracy at the $\alpha$ that gave the highest
accuracy for the same task and the same kind of set (same rule, same $N$) in fold~2, and the
reverse; the
choice of $\alpha$ is therefore never made on the items it is scored on. Random sets are scored
at the $\alpha$ chosen for the CC set of their size. Effects are averaged over the two folds.
The \emph{specific effect} is the task-mean accuracy minus the control-mean accuracy at the
same $\alpha$.

\paragraph{Statistics.} Two sets of the same task are patched on the same items, so they are
compared item by item. Each difference between two sets (on one task, or averaged over the 14
tasks) gets a 95\% interval from a paired bootstrap over the patched items within each task
(10{,}000 resamples, both folds resampled separately). The cross-fitted $\alpha$ is held fixed
in the resamples. A task counts for a set in the
per-task counts below only if the set's control effect is less than half of its effect on that
task.

\paragraph{Hypotheses.} All on the cross-fitted effect in check~1; the same comparisons in
check~2 are reported as secondary.
\begin{description}
  \item[H4 (lexicon).] CC has a larger effect than E1. Supported if the mean difference over
    the 14 tasks has a 95\% interval above 0 and CC is above E1 on at least 10 of the 14 tasks.
    Falsified if the interval contains 0 or CC is above E1 on 7 tasks or fewer.
  \item[H5 (not chance).] CC has a larger effect than random heads of its size and layers.
    Supported if its mean over the 14 tasks is above that of every random draw and, on at
    least 10 tasks, above the largest of the five draws. Falsified if its mean falls within the
    range of the five draws' means.
  \item[H6 (decoding against causal choice).] At each size $N \in \{20, 40, N_t\}$, CC-$N$
    is not worse than FV-$N$: the lower end of the 95\% interval of the mean difference
    (CC-$N$ minus FV-$N$) is above $-0.05$. In addition, CC-$N$ is above FV-$N$ on at least
    4 of the 5 tasks on which the E1 sets do nothing and the CC set did most in Sec.~15:
    antonym, next-item, English--French, English--German, English--Spanish. Decided
    separately for each $N$; falsified at a size if either part fails.
\end{description}
Expected from the earlier results, and stated so that the outcome can be compared with it:
H4 and H5 hold; H6 holds at $N_t$ and may fail at 20 and 40, where CC sets ranked by match
rate gave 0.21 and 0.25 (summary, E3) and the 30-head FV of the summary has a larger mean.

\paragraph{What each outcome would mean.}
\begin{itemize}
  \item H4 and H5 hold, H6 holds at every $N$: choosing heads by what they decode is as
    good as choosing them by their causal effect, and better on the tasks named in H6. This
    is the strongest outcome, and the one that would make Sec.~15 a main result.
  \item H4 and H5 hold, H6 only at $N_t$: decoded words find the right heads, but spread
    over more of them than a causal ranking needs; the claim becomes one about which heads,
    not how few.
  \item H4 holds, H5 fails: the CC sets work by their size and layers, not by which heads
    they hold; Sec.~15's gain over E1 is a size effect.
  \item H4 fails: Sec.~15's gain came from choosing the design on the test items.
\end{itemize}

\subsection{E4b: More wordings, and wordings held out}
\label{sec:e4b}

\paragraph{Question.} The five templates of Yang et al.\ are few, and E3's word lists and
heads were chosen on all of them, so nothing has been tested on a wording that was not used
to choose. Two questions. First, a head may carry the task but decode a list word only under
some wordings; then more wordings should find more of the heads that carry the task, as
pooling the five templates did in Sec.~15 (0.18 to 0.45, with five times the prompts of the
best template). Second, if a set holds the task's
representation and not the words of particular instructions, its heads should come up, and its
patch should work, under wordings never used to choose it, including wordings that do not
contain the task's name.

\paragraph{New wordings.} For each task, 30 new instruction templates, in six cells of five:
three forms $\times$ whether the instruction names the task.
\begin{itemize}
  \item Forms: \emph{cloze} (``The opposite of \{x\} is''), \emph{question--answer} (an
    instruction, then \verb|Q: {x}\nA:|, as Yang's templates B--E), and
    \emph{label} (two labels, one before the input and one after it, \verb|Word: {x}\nOpposite:|).
  \item Task named: the instruction contains a word of the task's term list of Yang et al.\ or
    the task's name (antonym, capital, French, \ldots). Not named: it contains none of them
    and describes the task otherwise (``the word whose meaning is furthest from'', ``the city where
    the government of \{x\} sits'', ``this word as a person in Lyon would say it''). Presence is
    checked by script on the lowercased words.
\end{itemize}
They are written by an LLM (GPT-5.5, with one fixed prompt per cell that gives the task's
description and two of Yang's templates as examples, not the items), checked by script for the
cell's conditions and by hand for correctness and ambiguity (one pass by the author), and the
first five per cell that pass are kept. A cell that cannot be filled (the not-named translation
cells may be hard) is reported with the number kept. Yang's five templates are kept as they are.

\paragraph{Behaviour (E0 on the new wordings).} Every new template is run on every item (at
most 500 per task, the same items for all templates), greedy, saved and scored as E0
(Sec.~4.5). Reported: accuracy per template, its spread per task under the flexible and the
first-token score, and the mean accuracy per form and per named / not named. These choose
$P^+$ below, on the choosing half only.

\paragraph{Wording split.} In each of the six cells, three templates are used to choose
(seed~0) and two are held out: 18 new templates to choose and 12 held out per task, the held-out
twelve two per cell. Yang's five templates are always among those used to choose, so that the
sets of E4 are the condition $m = 5$ below. Held-out wordings are never used to choose words,
heads, rankings or $\alpha$. The wording split is crossed with the item split: words and heads
are chosen on the choosing half under the chosen wordings, and every test is on the testing
half.

\paragraph{Conditions: number of wordings.} $m \in \{1, 5, 11, 23\}$ wordings to choose words
and heads, nested: $m = 1$, the most accurate of Yang's templates on the choosing half (the
``best template'' of Sec.~15); $m = 5$, Yang's five; $m = 11$, Yang's five and one new template
from each cell, three draws (seeds 0--2); $m = 23$, Yang's five and all 18 new ones. The number
of prompts is held fixed so that $m$ changes how varied the prompts are, not how many: each
condition uses $T_t$ correct prompts in all, where $T_t$ is the number of correct prompts of the
$m = 1$ template on the choosing half, capped at 500, split evenly over the $m$ templates (the
first correct prompts of each; a shortfall in one template is reported, not filled from
another). The $m = 5$ condition is therefore not the CC set of E4, which takes up to 100
prompts from each of the five templates: the two are the same set only where $T_t = 500$
(antonym, synonym and the three translations), and on the other tasks the $m = 5$ condition
has fewer prompts per template (about 15 to 75). The CC set of E4 is reported beside the
conditions as a reference point (five templates, up to 100 prompts each), not as one of them.
Secondary, at $m = 23$: 100 prompts per template, to see
whether more prompts add to more wordings.

\paragraph{Head sets.} For each $m$, the CC rule of E4 unchanged: word shares, task and
common words, and Definition~1 on the pooled prompts of the $m$ templates; the other 13
tasks' shares for the contrast are computed under the same $m$. Secondary: \textbf{E1-pooled},
Yang's term list with Definition~1 on the same pooled prompts, at $m = 5$ and $m = 23$, to see
whether more wordings help the hand-written lists. For the $m = 5$ condition, five random sets
with the same number of heads in each layer are drawn as in E4 (seeds 0--4), since the set
differs from E4's.

\paragraph{Tests.}
\begin{enumerate}
  \item \textbf{Number of wordings}: both checks of E4 for every set, with the same five means
    (Yang's templates, choosing half) and the same items, so that only the heads differ;
    cross-fitted $\alpha$ by condition ($m$). The CC set of E4 is shown beside the four
    conditions. Measures beside the effect: set size; overlap
    (Jaccard) between the three $m = 11$ draws and between consecutive $m$; overlap of the
    word lists.
  \item \textbf{Reach on held-out wordings}: for the E1 set, CC ($m = 5$) and CC ($m = 23$),
    the match rate of the set's heads (its own word list, Definition~1's per-prompt test) on the
    correct prompts of each held-out template on the choosing half, against five random sets
    with the same number of heads in each layer as the set (drawn as in E4); separately for
    named and not-named templates, with the copy share (Sec.~4.6) on each.
  \item \textbf{Held-out means}: check~1 on the same three sets, patched with the means of the
    12 held-out templates (each from its correct prompts on the choosing half) instead of
    Yang's five.
  \item \textbf{Held-out failures}: a second check on wordings, not shots. Patched prompts:
    the prompts of each held-out template answered wrongly on the testing half (the first 20
    per template, at most 240), whose unpatched accuracy is 0. Means: Yang's five, choosing
    half. Control: the control mean of E4. This asks whether heads and means taken from the
    chosen wordings repair the wordings on which the model fails, which is the prompt
    sensitivity of Yang et al.\ in its most direct form.
\end{enumerate}
Tests 3 and 4 are run at three values of $\alpha$: the cross-fitted one of the set and 4 and 8.

\paragraph{Hypotheses.} On the cross-fitted effect, mean over the 14 tasks, intervals by the
paired bootstrap of E4.
\begin{description}
  \item[H7 (more wordings).] (a) CC at $m = 5$ has a larger check-1 effect than at $m = 1$
    (interval above 0): pooling templates helps when the number of prompts is the same.
    (b) CC at $m = 23$ has a larger effect than at $m = 5$ (interval above 0). Each part decided
    separately. Expected: (a) holds but with a smaller gain than in Sec.~15, where pooling also
    gave five times the prompts; (b) is smaller and may fail, since the union rule (E1b, $k = 1$)
    helped only some tasks. Comparing $m = 5$ with the CC set of E4 shows what the extra prompts
    add at five templates; it is reported, not tested.
  \item[H8 (held-out wordings).] For CC at $m = 5$ and at $m = 23$, each decided separately:
    (a) the match rate on held-out templates is above that of every random draw, on the named
    and on the not-named templates; (b) check~1 with the held-out means is not worse than with
    Yang's means by more than 0.05 (lower end of the interval); (c) on held-out failures, the
    task mean's effect is above the control mean's (interval above 0). Falsified for a set if
    any part fails; parts are reported one by one.
  \item[H9 (copying).] For CC at $m = 5$, the match rate on not-named held-out templates is
    at least half of that on named ones. If it is less, the heads mostly decode list words
    that the instruction contains, and the reach of H8(a) comes from copying.
\end{description}

\paragraph{What the outcomes would mean.} H7(b) and H8 together decide whether ``more
wordings'' improves identification or only enlarges the sets: if (b) holds and the $m = 23$ set
also passes H8 where $m = 5$ fails, the extra wordings found heads that carry the task across
wordings. If H8 holds but H9 fails, the sets work across wordings that name the task and say
little about a representation of the task beyond its name. H8(c) is the result that speaks to
Yang et al.'s main claim: heads and means from some wordings repairing failures of others.

\paragraph{Cost of E4b.} New-wording E0: 30 templates $\times$ at most 500 items $\times$ 14
tasks $\le$ 210{,}000 replies. Readings (J-head top tokens of every head at the last token): the
23 chosen and 12 held-out templates' correct prompts, at most about 245{,}000 forward passes
without generation. Patching, per task and fold: test~1, 9 new sets (CC at $m = 1$, $m = 5$,
three at $m = 11$, $m = 23$, $m = 23$ with 100 per template, two E1-pooled) $\times$ 12{,}000
$\approx$ 108{,}000; tests~3 and 4, 3 sets $\times$ 3 values of $\alpha$ $\times$ (12 means
$\times$ 100 $+$ 5 means $\times$ 240 $+$ 240 control) $\approx$ 24{,}000; about 132{,}000 per
task and fold, 3.7 million in all.

\subsection{Cost and output}

\paragraph{Cost of E4.} Per task and fold, with $\le 100$ patched items and five task means:
12 main and secondary sets $\times$ 2 checks $\times$ (10 values of $\alpha$ $\times$ 5 means
$+$ 10 $\times$ 1 control) $\times$ 100 $\approx$ 144{,}000 replies; 15 random sets $\times$ 2
checks $\times$ 10 $\times$ 5 $\times$ 100 $\approx$ 150{,}000; in all about 294{,}000 per task
and fold, 8.2 million over 14 tasks and two folds, plus the function-vector scores on each
choosing half. Two cuts keep the hypotheses intact: random sets in check~2 only at the
cross-fitted $\alpha$ (saves about 67{,}000 per task and fold), and the five secondary sets in
fold~1 only (about 60{,}000 in fold~2). Replies can stop at 16 new tokens: the flexible score
reads only the start of the reply, and the longest reference has 15 tokens. E4 and E4b
together: about 11.9 million patched replies before the cuts, 0.2 million new E0 replies and
0.25 million reading passes. If only part can be run, E4 comes first: H4--H6 do not depend on
E4b, while E4b takes its item split, means and control from E4.

\paragraph{Output.} \texttt{results/heldout/fold1/} and \texttt{fold2/}, and for E4b
\texttt{results/heldout/wordings/}: the new templates with their cells and checks, the
wording split, the split (item ids per half), the
word lists and head sets of each fold, the FV scores, accuracies at every $\alpha$ for every set,
mean and check, the cross-fitted and best-$\alpha$ effects, and the bootstrap intervals. Tables:
cross-fitted effect per task and set in check~1 and check~2 (as Table~14), and the three
hypothesis comparisons with their intervals; figures: accuracy against $\alpha$ per task for
E1, CC, FV-$N_t$ and the random draws; for E4b, the effect against $m$ per task, and
match rate and effect on held-out templates by form and named / not named.
